\documentclass[../DoAn.tex]{subfiles}
\begin{document}

\begin{center}
    \Large{\textbf{ABSTRACT}}\\
\end{center}
\vspace{1cm}
Most decentralized exchanges (DEXs) today rely on automated market makers (AMMs), which price assets through a fixed mathematical formula based on the liquidity available in a pool. This design causes large trades to suffer noticeable slippage and exposes them to maximal extractable value (MEV) attacks such as sandwiching. Some protocols have moved toward intent-based trading, where users sign an order describing their desired trade and leave the search for a counterparty and the execution itself to third parties called fillers, as seen in UniswapX and CoW Protocol. However, a survey of these systems shows that most still settle orders on an all-or-nothing basis, provide little economic accountability for fillers who register but fail to deliver, and operate within a single chain only.

To address these gaps, this thesis adopts the intent-based trading model as its foundation and extends it with the mechanisms found missing in the surveyed systems, since this model separates the expression of trading intent from execution, lowering the cost for users and allowing execution to be optimized through competition among fillers.

Building on this direction, the thesis presents NeutronX, a system that combines partial-fill order matching through a Dutch-auction mechanism with EIP-712 structured signatures and Permit2, a dynamic staking-and-slashing scheme that holds fillers accountable, a fallback executor that guarantees an order is completed before its deadline, and a cross-chain extension that enables atomic swaps between two EVM-compatible chains through hash time-lock contracts with Merkle-tree-based partial fills.

The main contribution of this thesis is the design, implementation, and testing of a complete DEX aggregator on a simulated network environment, in which orders are filled partially and transparently, fillers are held accountable for their commitments, every order is guaranteed to execute before its deadline, and users can perform cross-chain swaps without relying on a custodial bridge.

\end{document}